\documentclass[a4paper,10pt]{article}

% Packages
\usepackage{shellesc}
\write18{mkdir -p build}
\pdfoutput=1

\usepackage{txfonts}
\usepackage{titlesec}
\usepackage[utf8]{inputenc}
\usepackage[T1]{fontenc}
\usepackage[french]{babel}
\usepackage{geometry}
\usepackage{graphicx}
\usepackage{multicol}
\usepackage{fancyhdr}
\usepackage{hyperref}
\usepackage{float}
\usepackage{setspace}
\usepackage{caption}
\usepackage{booktabs}
\usepackage{enumitem}

% MISE EN PAGE
\geometry{a4paper, top=15mm, bottom=15mm, textwidth=185mm}
\setlength{\columnsep}{6mm}
\setlength{\parindent}{0pt}

\titleformat{\section}
  {\normalfont\bfseries}
  {\thesection.}{0.5em}{}
\titlespacing{\section}{0pt}{5pt}{2pt}
\titleformat{\subsection}
  {\normalfont\itshape}
  {\thesubsection.}{0.5em}{}

\pagestyle{fancy}
\renewcommand{\footrulewidth}{0pt}
\fancyhf{}
\rhead{\includegraphics[scale=0.4]{Images/hepia_simple.png}}
\rfoot{\textit{\today}}
\cfoot{\thepage}

\captionsetup{font=footnotesize, labelsep=colon}

\setlist[itemize]{noitemsep, topsep=2pt, parsep=0pt, partopsep=0pt}

\begin{document}

% EN-TETE
\begin{center}
\vspace*{0mm}
{\large
Cahier des charges --- Mini-projet AMC 2026\\
\vspace{2mm}
\textbf{Voice Recorder} : enregistrement, transcription et résumé de discussions
}\\
\vspace{5mm}
Guilhem Rorier Vilardell\footnote{\textit{\href{mailto:guilhem.rorier-vilardell@hes-so.ch}{MT2, guilhem.rorier-vilardell@hes-so.ch}}}\\
\vspace{2mm}
\footnotesize{\textit{Haute école de paysagerie, d'ingénierie et d'architecture, Genève, Suisse}}\\
\vspace{3mm}
\end{center}

\rule{\linewidth}{0.4pt}

\begin{center}
\begin{tabular}{ll}
\textbf{Type de projet :} & Monôme \\
\textbf{Cours :}          & Application Mobile avec Capteurs (AMC) \\
\textbf{Année :}          & 2026 \\
\end{tabular}
\end{center}

\rule{\linewidth}{0.4pt}

\vspace{3mm}
\setlength{\parskip}{2pt}

% ============================================================
\section{Contexte et objectif}
% ============================================================

Application Python-Kivy permettant d'enregistrer une discussion via le \textbf{microphone} de l'appareil (\texttt{plyer.audio} sur Android, \texttt{speech\_recognition}/\texttt{pyaudio} sur Windows), de transcrire l'audio en texte, d'en générer un résumé et d'exporter le résultat en \texttt{.md}. Cas d'usage : prise de notes passive lors de réunions ou cours. L'application est fonctionnelle en \textbf{simulation Windows} et \textbf{déployable sur tablette Android} via Buildozer (\texttt{.apk}).

% ============================================================
\section{Fonctionnalités principales}
% ============================================================

\begin{itemize}
  \item \textbf{Enregistrement audio} : démarrage/arrêt via boutons, capture en \texttt{.wav} stocké temporairement.
  \item \textbf{Transcription} : envoi à Google Speech-to-Text via \texttt{speech\_recognition}, en français, connexion requise.
  \item \textbf{Résumé automatique} : appel à l'API \textbf{Google Gemini 2.5 Flash} (prompt structuré), clé API dans \texttt{config.env}.
  \item \textbf{Export Markdown} : transcription + résumé sauvegardés en \texttt{.md} et affichés dans l'interface.
  \item \textbf{Historique SQLite} : 10 derniers enregistrements (date, durée, extrait) consultables depuis un écran dédié.
  \item \textbf{Copier-coller} : bouton pour copier le \texttt{.md} affiché dans le presse-papier.
\end{itemize}

% ============================================================
\section{Fonctionnalités non abordées en cours}
% ============================================================

\begin{itemize}
  \item \textbf{API de transcription vocale} : \texttt{speech\_recognition} avec backend Google Speech-to-Text~\cite{speechrecognition}.
  \item \textbf{LLM via API REST} : intégration de Google Gemini pour la génération de résumés en langage naturel~\cite{geminiapi}.
  \item \textbf{SQLite embarqué} : historique persistant avec requêtes SQL dans une application Kivy~\cite{sqlite}.
  \item \textbf{Fichiers Markdown} : écriture, lecture et affichage de \texttt{.md} avec accès au système de fichiers Android.
\end{itemize}

% ============================================================
\section{Compatibilité et déploiement}
% ============================================================

\textbf{Windows} : exécution directe avec Python et micro du PC. \textbf{Android} : \texttt{.apk} généré via Buildozer~\cite{buildozer}, permissions \texttt{RECORD\_AUDIO}, \texttt{INTERNET} et \texttt{WRITE\_EXTERNAL\_STORAGE} déclarées dans \texttt{buildozer.spec}. La clé API est stockée dans \texttt{config.env} (exclu du dépôt via \texttt{.gitignore}).

% ============================================================
\section{Répartition des tâches}
% ============================================================

\begin{table}[H]
\centering
\caption{Liste des tâches du projet}
\label{tab:taches}
\footnotesize
\begin{tabular}{lp{9.5cm}}
\toprule
\textbf{Priorité} & \textbf{Tâche} \\
\midrule
1  & Architecture et interface graphique Kivy (écrans, navigation). \\
2  & Accès au microphone (\texttt{plyer.audio} Android, \texttt{speech\_recognition} Windows). \\
3  & Transcription via Google Speech-to-Text. \\
4  & Résumé via Google Gemini 2.5 Flash. \\
5  & Export et affichage du fichier \texttt{.md}. \\
6  & Base de données SQLite pour l'historique (10 entrées). \\
7  & Bouton de copie dans le presse-papier. \\
8  & Tests et débogage sur Windows (mode simulation). \\
9  & Génération du \texttt{.apk} via Buildozer et tests sur tablette Android. \\
10 & Rédaction du \texttt{README.md} et mise à jour du cahier des charges. \\
11 & GitEdu et \texttt{JOURNAL.md}. \\
12 & Vidéo de démonstration (SWITCHtube). \\
\bottomrule
\end{tabular}
\end{table}
\pagebreak
% ============================================================
\section{Références}
% ============================================================

\begin{thebibliography}{9}

\bibitem{speechrecognition}
SpeechRecognition Library,
\textit{SpeechRecognition 3.x --- Python speech recognition library},
GitHub, 2024.
\url{https://github.com/Uberi/speech_recognition}
(consulté le 29 avril 2026).

\bibitem{geminiapi}
Google,
\textit{Gemini API documentation --- Google AI Studio},
Google Developers, 2025.
\url{https://ai.google.dev/gemini-api/docs}
(consulté le 29 avril 2026).

\bibitem{sqlite}
SQLite Consortium,
\textit{SQLite Documentation},
2024.
\url{https://www.sqlite.org/docs.html}
(consulté le 29 avril 2026).

\bibitem{buildozer}
Kivy Organization,
\textit{Buildozer documentation --- Package a Kivy application for Android},
Read the Docs, 2024.
\url{https://buildozer.readthedocs.io/en/latest/}
(consulté le 29 avril 2026).

\end{thebibliography}

\end{document}
